\begin{examplebox}
\textbf{例 1：由定义求方程。}
$F_1(-4,0),F_2(4,0)$，$|PF_1-PF_2|=6$。
先查 $0<6<8$，再取 $a=3,c=4$，得 $b^2=7$。
故轨迹为 $\boxed{x^2/9-y^2/7=1}$。
回验：右顶点 $(3,0)$ 的两条焦半径是 $7,1$，差为 $6$。

\medskip
\textbf{例 2：更换焦点方向。}
$F_1(0,-3),F_2(0,3)$，距离差绝对值为 $4$。
$0<4<6$，$a=2,c=3,b^2=5$，故 $\boxed{y^2/4-x^2/5=1}$。
回验：$(0,2)$ 到两焦点距离为 $5,1$，差为 $4$。
\end{examplebox}
\begin{examplebox}
\textbf{例 3：两个正候选值也可能需要舍根。}
双曲线 $x^2/9-y^2/16=1$ 上有一点 $P$，到左焦点距离 $PF_1=7$，求 $PF_2$。
$a=3,c=5$，所以 $|7-PF_2|=6$，得到 $PF_2=1$ 或 $13$。
\begin{center}\begin{tikzpicture}[font=\small,>=stealth]
\draw[black!30,rounded corners=2pt] (0,0) rectangle (14.4,8.2);
\node[anchor=west,font=\small\bfseries] at (.5,7.72) {例 3：绝对值方程的两个正根，还要检查几何存在性};
\node at (3.6,6.65) {$PF_1=7,\ PF_2=1$};\begin{scope}[shift={(3.6,3.8)},scale=.145]
\draw[orange!85!black,thick] (-5,0) circle (7);
\draw[violet,thick] (5,0) circle (1);
\node[fill=black,circle,inner sep=1.5pt] at (-5,0) {};\node[below=4pt] at (-5,0) {$F_1$};
\node[fill=black,circle,inner sep=1.5pt] at (5,0) {};\node[below=4pt] at (5,0) {$F_2$};
\end{scope}
\node at (3.6,1.3) {两圆分离：不存在这样的点};
\node at (10.8,6.65) {$PF_1=7,\ PF_2=13$};\begin{scope}[shift={(10.8,3.8)},scale=.145]
\draw[orange!85!black,thick] (-5,0) circle (7);
\draw[violet,thick] (5,0) circle (13);
\node[fill=black,circle,inner sep=1.5pt] at (-5,0) {};\node[below=4pt] at (-5,0) {$F_1$};
\node[fill=black,circle,inner sep=1.5pt] at (5,0) {};\node[below=4pt] at (5,0) {$F_2$};
\node[fill=black,circle,inner sep=2pt] at (-6,6.9282032303) {};\node[above left=5pt] at (-6,6.9282032303) {$P$};
\end{scope}
\node at (10.8,1.3) {两圆相交：候选距离可实现};

\node[anchor=west] at (.7,.55) {同尺度；焦距 $10$。左图 $7+1<10$；右图交点同时满足两条距离要求};
\end{tikzpicture}\end{center}

若取 $1$，两距离之和 $8<10$，违反三角不等式，舍去。
若取 $13$，两圆确实相交；答案为 $\boxed{13}$。
\textbf{坐标回验：} 由距离平方差 $4cx=7^2-13^2=-120$ 得 $x=-6$；
代入双曲线得 $y=\pm4\sqrt3$。此时两距离确为 $7,13$。
不能只检查“两个根都为正”，也不能只说满足某一个三角不等式便结束存在性验证。
\end{examplebox}
\clearpage

% ===== C037 扩充区开始：原有 1--3 例与 TikZ 图保持不变 =====
\begin{examplebox}
\textbf{例 4：最不容易忘的数字示例——$3,4,5$。}
给定 $F_1(-5,0),F_2(5,0)$，距离差绝对值为 $6$。
\[
a=3,\quad c=5,\quad b=\sqrt{25-9}=4,
\quad\boxed{\frac{x^2}{9}-\frac{y^2}{16}=1}.
\]
右顶点 $A_2(3,0)$：$A_2F_1=8,A_2F_2=2$，差为 $6$；
左顶点 $A_1(-3,0)$：$A_1F_1=2,A_1F_2=8$，差的绝对值仍为 $6$。
\textbf{一句话识别：} 顶点数轴坐标是 $\pm3$，焦点是 $\pm5$，
$b=4$ 是方程参数而非纵轴截距。

\medskip
\textbf{例 5：看常数就判轨迹，不必动笔求方程。}
焦距 $F_1F_2=10$，依次要求 $|PF_1-PF_2|=K$：
\[
\begin{array}{c|c}
K=0&\text{垂直平分线}\\
K=6&\text{正常双曲线两支}\\
K=10&\text{两条外侧射线，含焦点}\\
K=12&\text{无解}
\end{array}
\]
原因始终是 $|PF_1-PF_2|\le10$，其中 $K=0$ 与 $K=10$
都不再是通常意义的双曲线。
\end{examplebox}
\begin{examplebox}
\textbf{例 6：考场直接求焦半径，并核对符号。}
在 $x^2/9-y^2/16=1$ 上取 $P=(6,4\sqrt3)$，
焦点为 $F_1(-5,0),F_2(5,0)$。
因为 $a=3,c=5,x=6\ge a$，右支速算
\[
PF_1=\frac{5\cdot6}{3}+3=13,\qquad
PF_2=\frac{5\cdot6}{3}-3=7.
\]
也可以直接用勾股核验：
\[
PF_1=\sqrt{11^2+(4\sqrt3)^2}=\sqrt{169}=13,
\quad PF_2=\sqrt{1^2+(4\sqrt3)^2}=\sqrt{49}=7.
\]
二者差 $6$，恰好等于 $2a$。
若把 $x$ 换成 $-6$，左右焦半径\textbf{交换数值}，不能仍将差写成正 $6$。

\medskip
\textbf{例 7：把中心平移，定义中的长度完全不变。}
焦点 $F_1(-3,-1),F_2(7,-1)$，距离差绝对值为 $6$。
中心为 $O(2,-1)$，所以 $c=5,a=3,b^2=16$。
焦点连线是水平直线，标准方程平移为
\[
\boxed{\frac{(x-2)^2}{9}-\frac{(y+1)^2}{16}=1}.
\]
\textbf{要点：} 先找中点，再看焦点沿哪个方向排布；
平移只是换原点，并没有改变焦距和两点之间的长度差。
\end{examplebox}
\clearpage
